% Literature-informed numerical benchmark. This does not update old numerical certificates.
\subsection{Parameter provenance and the revised numerical benchmark}
We distinguish material reference values, manufacturer specifications,
declared design choices, and quantities derived from the assembled model.
The benchmark is a planar, clamped--free, homogeneous solid circular
Euler--Bernoulli beam with design length $L=80$ mm and radius $r=0.6$ mm.
It is an idealized numerical model rather than an identified clinical device.
Gravity, external flow, contact, fluid added mass, and shear deformation
are excluded from this benchmark.

The material reference point is Sylgard 184 with approximately $10$ wt\%
curing agent, cured at $80\,^{\circ}\mathrm C$ for $2$ h.
Cho et al. report approximately $E=2.1$ MPa and
$\rho=1028\,\mathrm{kg\,m^{-3}}$ near this formulation
\cite{cho2021pdms}. These are approximate literature measurements, not
measurements of the present cylindrical specimen. The curing-agent fraction
corresponds to approximately $9{:}1$ base to curing agent by mass; it is not
identified with a $10{:}1$ formulation.

The magnetic component is a declared uncoated cubic-core design of edge
$a=0.250$ mm. N50 material specifications give
$B_r\in[1.40,1.45]$ T and density at least
$7500\,\mathrm{kg\,m^{-3}}$ \cite{smN50}.
We select $B_r=1.425$ T and $\rho_m=7500\,\mathrm{kg\,m^{-3}}$
as nominal modeling points. Using the engineering approximation
$\mu\simeq B_ra^3/\mu_0$ \cite{kjMoment}, with
$\mu_0\simeq4\pi10^{-7}\,\mathrm{H\,m^{-1}}$, gives
\[
 \mu=1.77184214\,10^{-5}\,\mathrm{A\,m^2},\qquad
 m_t=1.171875\,10^{-7}\,\mathrm{kg},\qquad
 J_t=1.220703125\,10^{-15}\,\mathrm{kg\,m^2}.
\]
These estimates describe the bare designed core. Product exterior dimensions,
coatings, adhesive mass, magnetization history, and specimen tolerances are
not identified with these nominal inputs. The core center of mass is
idealized at the beam endpoint. Its translational mass and transverse rotary
inertia are included in $M_{\rm total}$; an offset installation requires the
corresponding inertial coupling terms.

We retain the design bias $B_0=30$ mT. Published magnetic navigation
experiments report fields of up to $40$ mT in arbitrary directions for their
OctoMag system \cite{lussi2021cvs}. The combination of this bias with a
$20$ mT transverse component has magnitude $36.056$ mT. This is an
order-of-magnitude feasibility comparison, not a calibration of the present
field--gradient ports or their simultaneous reachable set. The minimal
closed-loop examples use algebraic ideal ports and report commanded fields;
they do not implement actuator saturation.

The linearized biased stiffness is
$K_b=K_{\rm beam}+\mu B_0s_Ls_L^{\mathsf T}$.
Consistent cubic-Hermite elements and mass-normalized eigenvectors give
\[
 K_b\Phi=M_{\rm total}\Phi\Omega^2,
 \qquad \Phi^{\mathsf T}M_{\rm total}\Phi=I.
\]
With $32$ elements, the first two frequencies are approximately
$(7.84000916,47.39957662)$ rad/s. The balanced two-mode coordinates use
\[
 \nu=(1,6.04585730),\quad
 d=(0.30722439,0.45042632)^{\mathsf T},\quad
 \ell=(0.0189673443,0.00775347614)^{\mathsf T}\,\mathrm m,
\]
and $E_0=\mu B_0=5.31552642\,10^{-7}$ J.
Mesh refinements and an independent continuous-beam endpoint characteristic
equation are supplied with the reproducibility files. This frequency check
does not establish nonlinear two-mode truncation accuracy.

No measured damping coefficient is asserted for this specimen and setup.
We declare the effective model $C=\alpha M_{\rm total}$ and the scenario
$\zeta_1=0.025$, referenced to the biased first frequency, so that
\[
 \alpha=2\zeta_1\omega_1=0.3920004581\,\mathrm{s^{-1}},
 \qquad \gamma=\alpha/\omega_1=0.05.
\]
The choices $\zeta_1\in\{0,0.01,0.025,0.05,0.1\}$ are sensitivity
scenarios rather than measured confidence limits. Material loss factors from
other Sylgard preparation and frequency conditions \cite{fischer2017adhesion}
are not substituted for modal damping ratios. Similarly, the synthetic
residual coefficient $\beta=0.2$ is an imposed algorithmic stress test.

For the nominal nonnegative physical energy, the clamped-beam inequality
$\|w\|_\infty^2\leq2L^3E/(3EI)$ gives
$E_c=3EIc^2/(2L^3)=10.0197171$ nJ at $c=4$ mm.
The conditional $0.95E_c$ budget is $9.5187312$ nJ.
This budget has not been imposed on the minimal control examples, whose
initial nominal energy already exceeds it; those runs therefore do not
establish the associated energy-based clearance guarantee.
Previously reported learned coefficients, COMSOL errors, interval KAM
certificates, and operational error budgets are not transferred to this
revised physical baseline without regeneration and verification.
